\clearpage
\section{Wybrane metody klasyfikacji}
\label{chap:metody}

\subsection{Zasady doboru i wspólna terminologia}
\label{sec:dobor-metod}

Metody wybrano ze względu na ich role w~porównaniu, a~nie wyłącznie datę publikacji lub dostępność kodu. Pula obejmuje siedem rozwiązań neuronowych i~cztery konwencjonalne. Pozwala ona porównać liniową klasyfikację, sąsiedztwo, zespoły drzew oraz odmienne sposoby wykorzystania sieci. Nie obejmuje wszystkich rodzin opisanych w~literaturze; brak metody w~macierzy nie oznacza jej nieskuteczności.

Przez \emph{potok uczenia} rozumie się tutaj przygotowanie danych, dopasowanie klasyfikatora i~regułę wyznaczania predykcji. Parametry modelu są wielkościami uzyskiwanymi w~trakcie dopasowania. Hiperparametry określają m.in. strukturę, regularyzację i~przebieg tego procesu, a~ich wybór stanowi osobny poziom selekcji. Konfiguracja domyślna jest wskazanym punktem odniesienia, lecz może zawierać ustawienia opracowane wcześniej na innych zadaniach.

Każda metoda ma zwracać prawdopodobieństwa w~jednakowej, zapisanej kolejności klas. Daje to wspólny interfejs oceny bez wymuszania identycznej reprezentacji wejścia. Narzucenie jednego kodowania wszystkim algorytmom mogłoby usunąć mechanizm, który uzasadnia ich wybór. Z~tego względu porównanie główne dotyczy kompletnych procedur. Tabela~\ref{tab:role-metod} podsumowuje ich role; szczegółowe ustawienia wykonania należą do metodyki i~artefaktów badania.

\begin{table}[tbp]
\centering
\caption{Metody przyjęte do porównania i ich role badawcze. Opracowanie własne; zestawienie nie jest rankingiem.}
\label{tab:role-metod}
\small
\begin{tabularx}{\textwidth}{@{}P{0.22\textwidth}P{0.24\textwidth}Y@{}}
\toprule
Metoda & Mechanizm & Rola w porównaniu \\
\midrule
Regresja logistyczna & Liniowa ocena klas & Odniesienie liniowe w przyjętej reprezentacji cech. \\
KNN & Sąsiedztwo przykładów & Konwencjonalne odniesienie do uczonego sąsiedztwa. \\
XGBoost & Wzmacniane drzewa & Silne odniesienie drzewiaste. \\
CatBoost & Drzewa i obsługa kategorii & Odniesienie zachowujące dedykowane przetwarzanie kategorii. \\
MLP & Warstwy gęste & Bazowa metoda neuronowa. \\
TabNet & Sekwencyjne maski cech & Specjalizacja przez selekcję informacji. \\
FT-Transformer & Tokeny cech i uwaga & Specjalizacja reprezentacji i interakcji kolumn. \\
RealMLP & Rozwinięta receptura MLP & Znaczenie przygotowania danych i uczenia. \\
ModernNCA & Uczona reprezentacja sąsiedztwa & Porównanie z klasycznym KNN. \\
TabM & Zespół ze współdzieleniem parametrów & Efektywne łączenie predykcji wewnątrz architektury. \\
TabPFN-3.5 & Pretrenowanie i kontekst zadania & Zewnętrznie nabyte wagi przy stałej adaptacji. \\
\bottomrule
\end{tabularx}
\end{table}

\subsection{Metody konwencjonalne}
\label{sec:metody-konwencjonalne}

\subsubsection{Regresja logistyczna}
\label{sec:regresja-logistyczna}

Regresja logistyczna wiąże liniową kombinację cech z~prawdopodobieństwem klasy. W~przypadku binarnym stosowana jest funkcja logistyczna, natomiast uogólnienie wieloklasowe pozwala uzyskać rozkład po etykietach. Regularyzacja ogranicza wartości współczynników. Takie własności opisuje również dokumentacja implementacji scikit-learn~\cite{sklearnLogistic}.

Model stanowi odniesienie pozwalające sprawdzić, czy do uzyskania dobrej jakości potrzebna jest bardziej złożona zależność niż liniowa w~przyjętej reprezentacji. Nie oznacza to liniowości względem surowego zapisu każdej kolumny: kodowanie kategorii zmienia wejście. W~przyjętym potoku nominalne wartości są kodowane zero-jedynkowo, a~przekształcenia dopasowywane na odpowiedniej części treningowej. Dobra jakość tego odniesienia ogranicza uzasadnienie kosztowniejszych metod; słaba jakość sama w~sobie nie dowodzi przewagi określonej architektury.

\subsubsection{Klasyfikacja najbliższych sąsiadów}
\label{sec:knn}

Reguła najbliższego sąsiada przypisuje nowemu przypadkowi klasę na podstawie podobnych oznakowanych obserwacji~\cite{cover1967nearest}. KNN rozszerza tę zasadę na wybrany zbiór sąsiadów i~agreguje ich etykiety, ewentualnie z~wagami zależnymi od odległości~\cite{sklearnNeighbors}. Kluczowe znaczenie mają liczba sąsiadów, metryka i~reprezentacja cech. Duża skala jednej kolumny może zdominować odległość, a~dowolne kody kategorii mogą narzucić nieuzasadnione sąsiedztwo.

W~badaniu nominalne kategorie KNN mają ustalone kodowanie zero-jedynkowe. Należy uwzględniać pamięć banku odniesienia i~koszt predykcji, a~nie tylko przygotowania klasyfikatora. Natywna walidacja LOO dla zadań IID wyłącza etykietę ocenianej obserwacji z~jej sąsiedztwa, lecz dopasowuje preprocessing na całej części deweloperskiej. Jest to jawny wyjątek walidacyjny, odmienny od dokładnego ponownego dopasowania wszystkich przekształceń po usunięciu każdego wiersza. Dla Diabetes zastosowane zostaną foldy grupowe.

\subsubsection{XGBoost}
\label{sec:xgboost}

XGBoost jest systemem wzmacniania drzew, w~którym predyktor powstaje przez kolejne dodawanie funkcji drzewiastych. Nowe drzewa ograniczają stratę aktualnego zespołu, a~kryterium uwzględnia regularyzację jego złożoności. Chen i~Guestrin~\cite{chen2016xgboost} opisują również mechanizmy usprawniające uczenie na danych rzadkich i~dużych kolekcjach.

Podziały drzew umożliwiają modelowanie progów i~interakcji bez przyjmowania globalnej liniowości. Nie usuwa to potrzeby doboru głębokości, tempa uczenia, liczby drzew i~regularyzacji. W~pracy XGBoost stanowi silne odniesienie konwencjonalne, oceniane z~odpowiednią funkcją klasyfikacji i~zadeklarowaną reprezentacją. Własności współczesnej biblioteki nie są utożsamiane z~każdym szczegółem publikacji z~2016 roku; wersja wykonania jest odrębnym elementem odtwarzalności.

\subsubsection{CatBoost}
\label{sec:catboost}

CatBoost łączy wzmacnianie drzew z~obsługą cech kategorycznych. Prokhorenkova i~współautorzy~\cite{prokhorenkova2018catboost} opisują statystyki kategorii wyznaczane według porządku obserwacji oraz uporządkowane wzmacnianie. Celem tych mechanizmów jest ograniczanie określonych źródeł przecieku i~przesunięcia predykcji związanych z~użyciem etykiet treningowych.

W~badaniu zostanie zachowana natywna obsługa cech kategorycznych CatBoost. Metoda będzie punktem odniesienia w~porównaniu głównym oraz w~badaniu ważenia klas. Efektywna konfiguracja określi, które warianty mechanizmów opisanych w~publikacji są używane przez bibliotekę.

\subsection{MLP i znaczenie receptury uczenia}
\label{sec:mlp}

Perceptron wielowarstwowy przetwarza wektor wejściowy przez kolejne warstwy gęste z~nieliniowymi aktywacjami. W~klasyfikacji końcowa warstwa wyznacza oceny klas przekształcane w~prawdopodobieństwa. Uczenie parametrów opiera się na minimalizacji straty z~wykorzystaniem gradientów. Prosty MLP jest istotnym odniesieniem, ponieważ pozwala ocenić, czy bardziej wyspecjalizowane mechanizmy dają korzyść względem odpowiednio skonfigurowanej sieci gęstej~\cite{rtdlImplementation}.

Lokalny model bazowy wykorzystuje aktywacje ReLU, dropout i~optymalizator AdamW. Konfiguracja domyślna zawiera dwie warstwy ukryte o~256 i~128 jednostkach, dropout równy~0,1, tempo uczenia~0,001 i~zanik wag~$10^{-4}$. Cechy numeryczne są skalowane, a~kategorie kodowane zero-jedynkowo z~limitem 256 kolumn na cechę. Uczenie trwa najwyżej 100 epok; zachowywany jest stan sieci o~najniższej stracie walidacyjnej. Są to ustawienia własnego modelu bazowego, a~nie dokładne odtworzenie jednej z~opublikowanych procedur MLP. W~wariancie strojonym część tych wartości jest wybierana z~zadeklarowanej przestrzeni hiperparametrów.

Dropout losowo wyłącza jednostki w~trakcie treningu, zmieniając reprezentacje dostępne w~kolejnych aktualizacjach~\cite{srivastava2014dropout}. Podczas zwykłej oceny jest wyłączony; procedura uwzględnia odpowiednie skalowanie. Zanik wag ogranicza parametry w~trakcie optymalizacji. W~AdamW jest oddzielony od aktualizacji wynikającej z~gradientu straty; nie należy utożsamiać go bez zastrzeżeń z~dodaniem kary $L_2$ do straty optymalizowanej Adamem~\cite{loshchilov2019adamw}.

Wpływ regularyzacji zależy od innych ustawień. Zbyt silne ograniczenie może prowadzić do niedopasowania. Niższy błąd treningowy nie musi oznaczać lepszej jakości na zbiorze testowym, a~mniejsza luka generalizacji może wynikać z~pogorszenia dopasowania treningu. Potrzebna jest wspólna analiza jakości na treningu i~na danych niewykorzystanych do aktualizacji parametrów.

W~osobnym badaniu zostanie zastosowana ustalona procedura uczenia MLP oparta na opublikowanym rozwiązaniu oraz dwie ablacje: wyłączenie dropoutu i~wyłączenie zaniku wag. Pozostałe ustawienia będą zachowane. Dobór tej receptury nie będzie oparty na wyniku strojenia z~pełną liczbą etykiet trzech badanych zadań. Wyniki wieloskładnikowych zestawów regularyzacji~\cite{kadra2021regularization} są motywacją takiego rozdzielenia, lecz nie określają automatycznie wartości parametrów własnej kontroli.

\begin{table}[tbp]
\centering
\caption{Dwa zastosowania MLP w pracy. Rozdzielenie modelu bazowego od kontrolowanego badania regularyzacji.}
\label{tab:mlp-role}
\small
\begin{tabularx}{\textwidth}{@{}P{0.24\textwidth}YY@{}}
\toprule
Właściwość & MLP w porównaniu głównym & MLP w badaniu regularyzacji \\
\midrule
Cel & Odniesienie neuronowe & Ocena skutku usunięcia składnika. \\
Dobór ustawień & Domyślne i strojenie wewnętrzne & Receptura ustalona niezależnie. \\
Liczność treningu & Wynik przyjętego podziału & Zagnieżdżone podzbiory treningu. \\
Kontrola zmian & Porównanie pełnych potoków & Jedna ablacja przy stałych pozostałych warunkach. \\
\bottomrule
\end{tabularx}
\end{table}

\subsection{TabNet i FT-Transformer}
\label{sec:tabnet-ft}

\subsubsection{TabNet: sekwencyjna selekcja cech}
\label{sec:tabnet}

TabNet wykorzystuje kolejne kroki przetwarzania, w~których uczone maski określają udział cech dla danej obserwacji. Informacja z~poprzedniego kroku wpływa na następną selekcję, a~częściowe reprezentacje współtworzą wynik. Mechanizm ma kierować zdolność modelowania na wybrane składniki wejścia zamiast przetwarzać je jednakowo w~każdym kroku~\cite{arik2021tabnet}.

Maski mogą służyć opisowi zachowania modelu, lecz nie są dowodem przyczynowego znaczenia cech. W~tej pracy TabNet reprezentuje historycznie ważną specjalizację dla tabel. Nie zaplanowano osobnego badania wierności wyjaśnień ani etapu samonadzorowanego pretrenowania, również opisanego w~oryginalnej publikacji. Wykorzystana implementacja biblioteczna będzie identyfikowana oddzielnie od publikacji metody.

\subsubsection{FT-Transformer: tokeny cech i uwaga}
\label{sec:ft-transformer}

FT-Transformer przekształca każdą cechę w~token o~wspólnym wymiarze. Dla cech numerycznych stosuje parametry właściwe kolumnie, a~dla kategorii --- uczone tablice osadzeń. Zachowuje to tożsamość cechy: jednakowa wartość liczbowa w~dwóch kolumnach nie musi otrzymać tej samej reprezentacji. Warstwy uwagi przetwarzają tokeny cech pojedynczej obserwacji, a~reprezentacja tokenu agregującego służy predykcji~\cite{gorishniy2021revisiting}.

Mechanizm pełnej uwagi wymaga rozpatrywania par tokenów, dlatego szerokość tabeli ma znaczenie dla kosztu. W~porównaniu z~TabNet inny jest sposób selekcji i~łączenia informacji. Różnica całych wyników nie wyizoluje jednak skutku samego rodzaju uwagi. Adapter w~repozytorium korzysta z~implementacji autorów; liczby bloków, wymiar tokenu i~ustawienia uczenia muszą być zapisane jako element wykonanej konfiguracji~\cite{rtdlImplementation}.

Na rysunku~\ref{fig:reprezentacje} przewidziano porównanie przepływu informacji w~MLP, TabNet i~FT-Transformer. Jest to schemat mechanizmów, a~nie ilustracja przewagi jednej metody.

\plannedfigure
    {figures/ch03/przeplyw-reprezentacji.pdf}
    {Przetwarzanie cech w MLP, TabNet i FT-Transformer. Opracowanie własne na podstawie opisów metod~\cite{gorishniy2021revisiting,arik2021tabnet}.}
    {fig:reprezentacje}
    {Trzy równoległe ścieżki od tej samej obserwacji do prawdopodobieństw klas. Pokazać wektor cech, kolejne maski lub tokeny z tożsamością kolumn.}

\subsection{RealMLP: rozwinięta receptura sieci gęstej}
\label{sec:realmlp}

RealMLP jest zestawem rozwiązań obejmujących przetwarzanie wejścia, konstrukcję sieci i~uczenie. W~publikacji opisano m.in. skalowanie odporne z~łagodnym ograniczaniem wartości, osadzenia i~zmodyfikowane ustawienia optymalizacji~\cite{holzmueller2024realmlp}. Nazwa nie powinna być interpretowana jako oznaczenie jednej dodatkowej warstwy dołączonej do dowolnego MLP.

W~badaniu RealMLP odpowiada na pytanie o~jakość rozwiniętego rozwiązania gęstego względem lokalnej sieci i~odniesień konwencjonalnych. Warianty domyślne i~strojone będą wynikać z~przypiętego adaptera i~jego przestrzeni konfiguracji. Konfiguracje do strojenia generuje biblioteka TabArena; lokalny plik YAML nie definiuje dla tej metody osobnej przestrzeni.

Ustawienia opracowane na innych zadaniach są formą wiedzy zewnętrznej. Wynik MLP--RealMLP będzie zatem różnicą kompletnych potoków, obejmującą również przygotowanie danych i~recepturę. Nie posłuży jako jednoczynnikowy dowód działania konkretnego regularyzatora; temu służy osobne badanie MLP.

\subsection{ModernNCA: uczenie reprezentacji sąsiedztwa}
\label{sec:modernnca}

ModernNCA rozwija ideę Neighbourhood Components Analysis, zastępując wyłącznie liniową reprezentację rozwiązaniem wykorzystującym współczesne składniki głębokiego uczenia. Przewidywanie wykorzystuje miękką agregację etykiet sąsiadów w~uczonej przestrzeni. Autorzy stosują również losowe próbkowanie kandydatów sąsiedztwa, aby ograniczać koszt treningu~\cite{ye2025modernnca}.

Sieć uczy się reprezentacji, w~której wyznaczane jest sąsiedztwo. Podczas predykcji nadal potrzebne są oznakowane przykłady odniesienia, dlatego koszt metody obejmuje również ich przechowywanie i~porównywanie z~nowymi obserwacjami.

KNN stanowi bezpośrednie odniesienie, lecz porównanie obejmie różne sposoby uczenia, przetwarzania i~agregacji. Nie wyizoluje wyłącznie korzyści reprezentacji. Kandydaci i~ich etykiety muszą pochodzić z~treningu bieżącego foldu; etykieta walidacyjna lub testowa nie może trafić do banku. W~środowisku użyto implementacji wywodzącej się z~TALENT i~udostępnionej przez TabArena, co wymaga odróżnienia algorytmu publikacji od konkretnej ścieżki integracji.

\subsection{TabM i poziomy łączenia predykcji}
\label{sec:tabm-zespoly}

TabM wykorzystuje sieć gęstą do wytwarzania wielu predykcji, przy czym podmodele współdzielą znaczną część parametrów. Rozwiązanie nawiązuje do BatchEnsemble i~ogranicza konieczność przechowywania całkowicie niezależnej kopii każdej sieci. W~klasyfikacji predykcje podmodeli są łączone w~ocenę klas~\cite{gorishniy2025tabm}.

Oszczędność parametrów nie oznacza zerowego kosztu dodatkowych predykcji. Liczba podmodeli, sposób ich uczenia i~agregacji należą do definicji wykonanego wariantu. W~tej pracy TabM reprezentuje efektywne zespoły neuronowe. Samo zestawienie go z~lokalnym MLP nie jest kontrolą, w~której jedyną różnicą jest uśrednianie: zmienia się również parametryzacja i~procedura uczenia.

Należy odróżnić trzy poziomy łączenia, zestawione w~tabeli~\ref{tab:poziomy-zespolow}. Wewnętrzny zespół TabM występuje wewnątrz pojedynczego dopasowania modelu. Uśrednienie foldów łączy modele uczone na różnych częściach danych. Zespół konfiguracji łączy z~kolei wyniki różnych kandydatów, z~wagami ustalonymi na walidacji. Poziomy mogą się nakładać; nie są wzajemnie wykluczającymi się warunkami.

\begin{table}[tbp]
\centering
\caption{Poziomy łączenia predykcji. Opis definicji metod i przyjętego protokołu, bez własnych wyników jakości.}
\label{tab:poziomy-zespolow}
\small
\begin{tabularx}{\textwidth}{@{}P{0.22\textwidth}YY@{}}
\toprule
Poziom & Łączone obiekty & Znaczenie dla badania \\
\midrule
Wewnętrzny & Podmodele ze współdzieleniem parametrów & Składnik architektury TabM~\cite{gorishniy2025tabm}. \\
Foldy & Modele tej samej konfiguracji & Zwykłe warunki D i T już wykorzystują takie uśrednienie. \\
Konfiguracje & Wyniki różnych kandydatów & T+E dodaje dobór wag na OOF, bez selekcji na teście. \\
\bottomrule
\end{tabularx}
\end{table}

Dla metod strojonych z~ośmiofoldową walidacją konfiguracja domyślna (D) i~wybrana konfiguracja (T) korzystają z~uśrednienia predykcji ośmiu foldów. Warunek T+E dodatkowo łączy konfiguracje. Nie są to trzy niezależne kampanie treningowe ani porównanie jednej sieci z~ośmioma. Natywny IID KNN i~stały TabPFN mają odmienne zasady, które należy zachować w~opisie wyników. Rysunek~\ref{fig:zespoly} ma uwidocznić właśnie to rozróżnienie.

\plannedfigure
    {figures/ch03/poziomy-zespolow.pdf}
    {Wewnętrzny zespół, uśrednienie foldów i łączenie konfiguracji. Opracowanie własne na podstawie mechanizmu TabM~\cite{gorishniy2025tabm} i przyjętego protokołu.}
    {fig:zespoly}
    {Trzy panele z podpisanymi jednostkami łączenia. Zaznaczyć, że D i T zwykłych metod już uśredniają foldy, a KNN i TabPFN mają wyjątki.}

\subsection{TabPFN-3.5: pretrenowanie i kontekst zadania}
\label{sec:tabpfn}

TabPFN-3.5 wykorzystuje wagi uzyskane przed otrzymaniem konkretnego zadania. Oznakowane obserwacje uczące stanowią kontekst, na podstawie którego model przewiduje etykiety nowych przypadków. Lokalna adaptacja nie jest utożsamiana z~trenowaniem całej sieci od początku. Raport techniczny~\cite{jaeger2026tabpfn35} rozróżnia model bazowy oraz warianty o~innych możliwościach i~kosztach inferencji.

W~pracy przyjęto jeden lokalny checkpoint bazowy, ze stałymi ustawieniami inferencji, bez HPO i~bez ośmiofoldowego zespołu treningowego. Kontekst pochodzi z~pełnej części deweloperskiej. Wewnętrzne ustawienie liczby estymatorów jest elementem tej stałej procedury i~nie tworzy warunku T+E. Warianty Plus lub Thinking nie są zamiennikami wybranego checkpointu.

Porównanie kosztu dotyczyć będzie lokalnego przygotowania i~predykcji, bez historycznego kosztu pretrenowania. Wcześniejsze publikacje TabPFN nie potwierdzają automatycznie możliwości ani wyników wersji 3.5. Ponadto zbiorcze zestawienia rodziny w~raporcie technicznym nie muszą dotyczyć jednego wariantu bazowego~\cite{jaeger2026tabpfn35}. Faktyczna dostępność checkpointu, obsługa rozmiarów i~klas oraz wymagania pamięciowe wymagają kwalifikacji. Ich opis nie będzie zastąpiony wynikami autorów raportu lub testem samego interfejsu.

\subsection{Granice porównań i uzasadnienie wyłączeń}
\label{sec:granice-metod}

Przyjęta pula obejmuje odmienne mechanizmy, ale nie całość literatury. ResNet, NODE, podejścia samonadzorowane, inne modele uwagi i~kolejne modele fundamentowe pozostają kontekstem dla wyboru. ResNet jest ważnym odniesieniem historycznym, lecz rozwinięty model gęsty reprezentuje w~macierzy RealMLP. ModernNCA reprezentuje uczenie sąsiedztwa. FT-Transformer wykorzystuje uwagę nad tokenami cech, a~TabM --- współdzielenie parametrów zespołu. Dodawanie kolejnych podobnych ról zwiększałoby koszt bez automatycznego pogłębienia odpowiedzi na przyjęte pytania.

Tak ograniczona pula pozwala zbadać wybrane sposoby wykorzystania sieci przy dostępnym budżecie obliczeń. Wnioski z~porównania będą dotyczyły tych konkretnych procedur i~warunków oceny. Wpływ ważenia klas i~regularyzacji zostanie oceniony osobno, zgodnie z~projektem przedstawionym w~kolejnym rozdziale.